\documentclass[titlepage,a4paper]{jsarticle}
\usepackage{../../../sty/import}% 各種パッケージインポート
\usepackage{../../../sty/title}% タイトルページの変更

%% タイトルページの変数
% レポートタイトル
\title{原子力材料と核燃料レポート}
% 提出日
\expdate{\today}
% 科目名
\subject{原子力材料と核燃料}
% 分野
\class{情報経営システム工学分野}
% 学年
\grade{M1}
% 学籍番号
\mynumber{24336488}
% 記述者
\author{本間三暉}

\begin{document}
% titleページ作成
\maketitle
\section{ステンレス鋼の分類を示せ．}
ステンレス鋼は主に$Cr$系ステンレス鋼，$Cr-Ni$系ステンレス鋼が存在する．

$Cr$系ステンレス鋼は，$Cr$を主成分とするステンレス鋼である．

$Cr$系ステンレス鋼には，マルテンサイト系ステンレス鋼，フェライト系ステンレス鋼がある．

マルテンサイト系ステンレスは，熱処理(主に焼入れ焼戻し)によって強さやじん性を付加することができるため，
耐摩耗性と同時に耐食性も重要な医科用機械器具，刃物，プラスチック金型などに使用されている．
13\%Crを主体としている(SUS410)と18\%Crを主体としている(SUS416)があり，後者のほうが耐久性には優れている．

フェライト系ステンレスは焼入れによって硬化しないもので，Cr含有量によって，低$Cr$系(11～14\%$Cr$)、中$Cr$系(16～18\%$Cr$)、高$Cr$系(25\%以上$Cr$)に分類でき，
$Ni$は含有しない．低$Cr$系や中$Cr$系は耐食性が低く，腐食性の比較的弱い環境で使用され，特に大気中での耐食性が良好であるため，建築物や厨房関係に用いられる．
フェライト相という体心立方晶(bcc)を構造に持つ相からなり，17\%Crが基本である(SUS430)．

$Cr-Ni$系ステンレス鋼は，$Cr$と$Ni$を主成分とするステンレス鋼である．

$Cr-Ni$系ステンレス鋼には，オーステナイト系ステンレス鋼と2層ステンレス，析出硬化型ステンレスがある．

オーステナイト系ステンレス鋼は，$Cr$の他に数\%の$Ni$を含有しており，全般的には他の系よりも耐食性に優れている．
オーステナイト系ステンレス鋼の代表的なものは18\%Crと8\%Ni(SUS304)である．

2層ステンレスは，オーステナイト相とフェライト相よりからなる，フェライト安定化元素の$Cr$がオーステナイト系より多く，$Cr/ Ni$がオーステナイト系より大きい．
オーステナイト系ステンレス鋼ではよく問題になる，塩素イオンによる応力腐食割れに対して優れた特性を持つため，海水や工業用水を用いるポンプなどに利用される

析出硬化型ステンレスは，金属間化合物による析出硬化で耐食性と高強度を兼ね備えており，SUS631は18\%Crと8\%Niの析出硬化型である．
SUS631は板材や線材として強さやばね特性を要する部品に使用されている．

% 参考文献
\begin{thebibliography}{99}
\bibitem{} 原子力材料と核燃料 第三回講義資料
\bibitem{} ステンレス鋼の種類と用途 \url{https://www.monotaro.com/note/readingseries/kikaibuhinhyomensyori/0401/}
\end{thebibliography}

\end{document}